\chapter{Stato dell'arte e fondamenti}
\label{chap:state-of-the-art}

% Apertura: perimetro e criteri di selezione della letteratura.
% Organizzare per problemi, evitando una divisione rigida fra metodi
% algoritmici e RL che produca sovrapposizioni.
% Bibliografia ragionata: thesis_latex/ricerca_bibliografica/.
% Il survey locale in fase di scrittura non deve essere citato.

\section{Navigazione autonoma dei robot a zampe}
\label{sec:sota-navigazione}

\subsection{Architetture gerarchiche di navigazione e locomozione}
% Distinguere pianificazione, comandi di navigazione e controllo locomotorio.

\subsection{Approcci integrati ed end-to-end}
% Confrontare architetture integrate e gerarchiche, delimitando il livello
% al quale interviene il contributo della tesi.

\section{Traversabilit\`a e rappresentazione del terreno}
\label{sec:sota-rappresentazioni}
% Quali informazioni sull'ambiente possono sostenere la navigazione?

\subsection{Rappresentazioni geometriche e mappe di elevazione}
% Geometria locale, pendenza, rugosita e rappresentazioni spaziali del terreno.

\subsection{Rappresentazioni semantiche e classificazione dei terreni}
% Ruolo e limiti delle categorie di terreno nell'assegnazione dei costi.

\subsection{Feature visive apprese e rappresentazioni senza classi esplicite}
% Modelli preaddestrati, DINO e DINOv2, limitatamente alle proprieta utili
% alla tesi. L'impiego concreto di feature, PCA e gridmap e nel capitolo 3.

\section{Apprendimento delle propriet\`a di attraversamento dall'esperienza}
\label{sec:sota-esperienza}
% Come collegare cio che il robot vede alle conseguenze del movimento?

\subsection{Informazioni esterocettive e propriocettive}
% Complementarita dei segnali visivi, geometrici, di contatto e di movimento.

\subsection{Apprendimento auto-supervisionato dall'interazione robot--terreno}
% Associazione fra osservazioni e misure raccolte durante il moto;
% distinguere costo, traversabilita, proprieta fisiche e stabilita.

\subsection{Memoria delle osservazioni e adattamento durante la navigazione}
% Distinguere memoria spaziale, storia propriocettiva e aggiornamento online
% dei parametri di un modello o di una policy.

\section{Efficienza energetica nei robot a zampe}
\label{sec:sota-energia}

\subsection{Grandezze energetiche e indicatori di efficienza}
% Consumo elettrico, lavoro meccanico e indicatori indiretti.
% La misura effettivamente implementata e descritta nel capitolo 3.

\subsection{Ottimizzazione energetica della locomozione}
% Discussione mirata su andature e movimenti articolari energeticamente
% efficienti, chiarendo la diversa scala di controllo rispetto alla tesi.

\subsection{Navigazione e pianificazione con costi energetici}
% Dedicare maggiore spazio alla scelta di tragitto e velocita in funzione
% del costo di attraversamento e ai lavori piu vicini al progetto.

\section{Reinforcement learning per la navigazione}
\label{sec:sota-rl}

\subsection{Formulazione del problema decisionale e osservabilit\`a parziale}
% Osservazioni, azioni, ritorno e funzioni di valore: introdurre soltanto
% gli strumenti necessari per comprendere il metodo.

\subsection{Apprendimento actor-critic: principi di PPO e SAC}
% Principi degli algoritmi, distinzione on-policy/off-policy, GAE e massima
% entropia quanto necessario. Configurazioni e reward sono nel capitolo 3.

\subsection{Addestramento in simulazione e uso di informazioni privilegiate}
% Simulazione parallela, apprendimento asimmetrico e differenza fra
% informazioni disponibili durante training ed esecuzione.

\section{Confronto critico e collocazione del lavoro di tesi}
\label{sec:sota-confronto}
% Tabella dei lavori vicini: sensori, rappresentazione, propriocezione,
% classi esplicite, costo, livello di controllo, memoria/adattamento e
% validazione. Motivare le scelte senza dedurre proprieta non verificate.
